\IfFileExists{stacks-project.cls}{%
\documentclass{stacks-project}
}{%
\documentclass{amsart}
}

% For dealing with references we use the comment environment
\usepackage{verbatim}
\newenvironment{reference}{\comment}{\endcomment}
%\newenvironment{reference}{}{}
\newenvironment{slogan}{\comment}{\endcomment}
\newenvironment{history}{\comment}{\endcomment}

% For commutative diagrams we use Xy-pic
\usepackage[all]{xy}

% We use 2cell for 2-commutative diagrams.
\xyoption{2cell}
\UseAllTwocells

% We use multicol for the list of chapters between chapters
\usepackage{multicol}

% This is generally recommended for better output
\usepackage{lmodern}
\usepackage[T1]{fontenc}

% For cross-file-references
\usepackage{xr-hyper}

% Package for hypertext links:
\usepackage{hyperref}

% For any local file, say "hello.tex" you want to link to please
% use \externaldocument[hello-]{hello}
\externaldocument[introduction-]{introduction}
\externaldocument[conventions-]{conventions}
\externaldocument[sets-]{sets}
\externaldocument[categories-]{categories}
\externaldocument[topology-]{topology}
\externaldocument[sheaves-]{sheaves}
\externaldocument[sites-]{sites}
\externaldocument[stacks-]{stacks}
\externaldocument[fields-]{fields}
\externaldocument[algebra-]{algebra}
\externaldocument[brauer-]{brauer}
\externaldocument[homology-]{homology}
\externaldocument[derived-]{derived}
\externaldocument[simplicial-]{simplicial}
\externaldocument[more-algebra-]{more-algebra}
\externaldocument[smoothing-]{smoothing}
\externaldocument[modules-]{modules}
\externaldocument[sites-modules-]{sites-modules}
\externaldocument[injectives-]{injectives}
\externaldocument[cohomology-]{cohomology}
\externaldocument[sites-cohomology-]{sites-cohomology}
\externaldocument[dga-]{dga}
\externaldocument[dpa-]{dpa}
\externaldocument[sdga-]{sdga}
\externaldocument[hypercovering-]{hypercovering}
\externaldocument[schemes-]{schemes}
\externaldocument[constructions-]{constructions}
\externaldocument[properties-]{properties}
\externaldocument[morphisms-]{morphisms}
\externaldocument[coherent-]{coherent}
\externaldocument[divisors-]{divisors}
\externaldocument[limits-]{limits}
\externaldocument[varieties-]{varieties}
\externaldocument[topologies-]{topologies}
\externaldocument[descent-]{descent}
\externaldocument[perfect-]{perfect}
\externaldocument[more-morphisms-]{more-morphisms}
\externaldocument[flat-]{flat}
\externaldocument[groupoids-]{groupoids}
\externaldocument[more-groupoids-]{more-groupoids}
\externaldocument[etale-]{etale}
\externaldocument[chow-]{chow}
\externaldocument[intersection-]{intersection}
\externaldocument[pic-]{pic}
\externaldocument[weil-]{weil}
\externaldocument[adequate-]{adequate}
\externaldocument[dualizing-]{dualizing}
\externaldocument[duality-]{duality}
\externaldocument[discriminant-]{discriminant}
\externaldocument[derham-]{derham}
\externaldocument[local-cohomology-]{local-cohomology}
\externaldocument[algebraization-]{algebraization}
\externaldocument[curves-]{curves}
\externaldocument[resolve-]{resolve}
\externaldocument[models-]{models}
\externaldocument[functors-]{functors}
\externaldocument[equiv-]{equiv}
\externaldocument[pione-]{pione}
\externaldocument[etale-cohomology-]{etale-cohomology}
\externaldocument[proetale-]{proetale}
\externaldocument[relative-cycles-]{relative-cycles}
\externaldocument[more-etale-]{more-etale}
\externaldocument[trace-]{trace}
\externaldocument[crystalline-]{crystalline}
\externaldocument[spaces-]{spaces}
\externaldocument[spaces-properties-]{spaces-properties}
\externaldocument[spaces-morphisms-]{spaces-morphisms}
\externaldocument[decent-spaces-]{decent-spaces}
\externaldocument[spaces-cohomology-]{spaces-cohomology}
\externaldocument[spaces-limits-]{spaces-limits}
\externaldocument[spaces-divisors-]{spaces-divisors}
\externaldocument[spaces-over-fields-]{spaces-over-fields}
\externaldocument[spaces-topologies-]{spaces-topologies}
\externaldocument[spaces-descent-]{spaces-descent}
\externaldocument[spaces-perfect-]{spaces-perfect}
\externaldocument[spaces-more-morphisms-]{spaces-more-morphisms}
\externaldocument[spaces-flat-]{spaces-flat}
\externaldocument[spaces-groupoids-]{spaces-groupoids}
\externaldocument[spaces-more-groupoids-]{spaces-more-groupoids}
\externaldocument[bootstrap-]{bootstrap}
\externaldocument[spaces-pushouts-]{spaces-pushouts}
\externaldocument[spaces-chow-]{spaces-chow}
\externaldocument[groupoids-quotients-]{groupoids-quotients}
\externaldocument[spaces-more-cohomology-]{spaces-more-cohomology}
\externaldocument[spaces-simplicial-]{spaces-simplicial}
\externaldocument[spaces-duality-]{spaces-duality}
\externaldocument[formal-spaces-]{formal-spaces}
\externaldocument[restricted-]{restricted}
\externaldocument[spaces-resolve-]{spaces-resolve}
\externaldocument[formal-defos-]{formal-defos}
\externaldocument[defos-]{defos}
\externaldocument[cotangent-]{cotangent}
\externaldocument[examples-defos-]{examples-defos}
\externaldocument[algebraic-]{algebraic}
\externaldocument[examples-stacks-]{examples-stacks}
\externaldocument[stacks-sheaves-]{stacks-sheaves}
\externaldocument[criteria-]{criteria}
\externaldocument[artin-]{artin}
\externaldocument[quot-]{quot}
\externaldocument[stacks-properties-]{stacks-properties}
\externaldocument[stacks-morphisms-]{stacks-morphisms}
\externaldocument[stacks-limits-]{stacks-limits}
\externaldocument[stacks-cohomology-]{stacks-cohomology}
\externaldocument[stacks-perfect-]{stacks-perfect}
\externaldocument[stacks-introduction-]{stacks-introduction}
\externaldocument[stacks-more-morphisms-]{stacks-more-morphisms}
\externaldocument[stacks-geometry-]{stacks-geometry}
\externaldocument[moduli-]{moduli}
\externaldocument[moduli-curves-]{moduli-curves}
\externaldocument[examples-]{examples}
\externaldocument[exercises-]{exercises}
\externaldocument[guide-]{guide}
\externaldocument[desirables-]{desirables}
\externaldocument[coding-]{coding}
\externaldocument[obsolete-]{obsolete}
\externaldocument[fdl-]{fdl}
\externaldocument[index-]{index}

% Theorem environments.
%
\theoremstyle{plain}
\newtheorem{theorem}[subsection]{Theorem}
\newtheorem{proposition}[subsection]{Proposition}
\newtheorem{lemma}[subsection]{Lemma}

\theoremstyle{definition}
\newtheorem{definition}[subsection]{Definition}
\newtheorem{example}[subsection]{Example}
\newtheorem{exercise}[subsection]{Exercise}
\newtheorem{situation}[subsection]{Situation}

\theoremstyle{remark}
\newtheorem{remark}[subsection]{Remark}
\newtheorem{remarks}[subsection]{Remarks}

\numberwithin{equation}{subsection}

% Macros
%
\def\lim{\mathop{\mathrm{lim}}\nolimits}
\def\colim{\mathop{\mathrm{colim}}\nolimits}
\def\Spec{\mathop{\mathrm{Spec}}}
\def\Hom{\mathop{\mathrm{Hom}}\nolimits}
\def\Ext{\mathop{\mathrm{Ext}}\nolimits}
\def\SheafHom{\mathop{\mathcal{H}\!\mathit{om}}\nolimits}
\def\SheafExt{\mathop{\mathcal{E}\!\mathit{xt}}\nolimits}
\def\Sch{\mathit{Sch}}
\def\Mor{\mathop{\mathrm{Mor}}\nolimits}
\def\Ob{\mathop{\mathrm{Ob}}\nolimits}
\def\Sh{\mathop{\mathit{Sh}}\nolimits}
\def\NL{\mathop{N\!L}\nolimits}
\def\CH{\mathop{\mathrm{CH}}\nolimits}
\def\proetale{{pro\text{-}\acute{e}tale}}
\def\etale{{\acute{e}tale}}
\def\QCoh{\mathit{QCoh}}
\def\Ker{\mathop{\mathrm{Ker}}}
\def\Im{\mathop{\mathrm{Im}}}
\def\Coker{\mathop{\mathrm{Coker}}}
\def\Coim{\mathop{\mathrm{Coim}}}

% Boxtimes
%
\DeclareMathSymbol{\boxtimes}{\mathbin}{AMSa}{"02}

%
% Macros for moduli stacks/spaces
%
\def\QCohstack{\mathcal{QC}\!\mathit{oh}}
\def\Cohstack{\mathcal{C}\!\mathit{oh}}
\def\Spacesstack{\mathcal{S}\!\mathit{paces}}
\def\Quotfunctor{\mathrm{Quot}}
\def\Hilbfunctor{\mathrm{Hilb}}
\def\Curvesstack{\mathcal{C}\!\mathit{urves}}
\def\Polarizedstack{\mathcal{P}\!\mathit{olarized}}
\def\Complexesstack{\mathcal{C}\!\mathit{omplexes}}
% \Pic is the operator that assigns to X its picard group, usage \Pic(X)
% \Picardstack_{X/B} denotes the Picard stack of X over B
% \Picardfunctor_{X/B} denotes the Picard functor of X over B
\def\Pic{\mathop{\mathrm{Pic}}\nolimits}
\def\Picardstack{\mathcal{P}\!\mathit{ic}}
\def\Picardfunctor{\mathrm{Pic}}
\def\Deformationcategory{\mathcal{D}\!\mathit{ef}}
\usepackage{cleveref}
\crefname{originallemma}{original lemma}{original lemmas}
\Crefname{originallemma}{Original lemma}{Original lemmas}
\makeatletter
\newlabel{lemma-third-isomorphism-triangle}{{014A}{source}{}{sourcepointer-1}{}}
\newlabel{lemma-third-isomorphism-triangle@cref}{{[originallemma][0][]014A}{[1][0][]source}{}{}{}}
\makeatother
\crefname{originallemma}{original lemma}{original lemmas}
\Crefname{originallemma}{Original lemma}{Original lemmas}
\makeatletter
\newlabel{lemma-third-map-square-zero}{{05QP}{source}{}{sourcepointer-2}{}}
\newlabel{lemma-third-map-square-zero@cref}{{[originallemma][0][]05QP}{[1][0][]source}{}{}{}}
\makeatother
\title{The homotopy colimit of a subsequence maps isomorphically to the homotopy colimit}
\author{The Stacks Project Authors}
\date{stacks-project:derived.tex:local-snapshot-20261012; GNU FDL 1.2 or later}
\begin{document}
\maketitle
\noindent\textit{Editorial note.} Counted span is the author's own proof environment, source lines 10448-10514, following the statement at 10430-10446.\par\medskip
\noindent\textit{Editorial note.} The h identities are 'the reader verifies' (10507): mechanical index bookkeeping from the displayed component formulas; the square-zero step is a '(small argument omitted)' instantiation of lemma-third-map-square-zero, recorded.\par\medskip
\noindent\textit{Editorial note.} Steps the author leaves to the reader are flagged here for the reviewer: source line(s) 10510 carry an "omitted"/"follows easily"/"immediate from the definitions" phrase. Those sentences are retained verbatim; the displayed argument around them stays inside the counted span.\par\medskip
\noindent\textit{Allowed standard prerequisites (editorial).}\begin{itemize}
\item Homotopic maps agree after localization. lemma-third-isomorphism-triangle, cited at proof line 10496; complete proof at source lines 641-670, queued for separate carding. Supplies that c a and d b being identities makes the induced map on homotopy colimits an isomorphism one way.
\item Nilpotent perturbations of the identity. lemma-third-map-square-zero, cited at proof line 10510; complete proof at source lines 710-728, queued for separate carding. Supplies that id - psi phi having square zero forces psi phi (hence phi, psi) to be isomorphisms.
\end{itemize}\medskip
\ifdefined\hypertarget\hypertarget{sourcepointer-1}{}\fi
\noindent\textit{Source reference.} Original lemma 014A: sources/stacks-derived/pinned-original-source.tex, source lines 631--631. Navigation pointer only: the counted proof cites this statement by \textbackslash{}ref as a standing input; the deduction the counted proof has to make is carried out inside the counted span, and no part of that deduction is delegated to this pointer.\par\medskip
\ifdefined\hypertarget\hypertarget{sourcepointer-2}{}\fi
\noindent\textit{Source reference.} Original lemma 05QP: sources/stacks-derived/pinned-original-source.tex, source lines 701--701. Navigation pointer only: the counted proof cites this statement by \textbackslash{}ref as a standing input; the deduction the counted proof has to make is carried out inside the counted span, and no part of that deduction is delegated to this pointer.\par\medskip
\begin{lemma}
\label{lemma-hocolim-subsequence}
Let $\mathcal{D}$ be a triangulated category.
Let $(K_n, f_n)$ be a system of objects of $\mathcal{D}$.
Let $n_1 < n_2 < n_3 < \ldots$ be a sequence of integers.
Assume $\bigoplus K_n$ and $\bigoplus K_{n_i}$ exist.
Then there exists an isomorphism
$\text{hocolim} K_{n_i} \to \text{hocolim} K_n$
such that
$$
\xymatrix{
K_{n_i} \ar[r] \ar[d]_{\text{id}} & \text{hocolim} K_{n_i} \ar[d] \\
K_{n_i} \ar[r] & \text{hocolim} K_n
}
$$
commutes for all $i$.
\end{lemma}

\begin{proof}
Let $g_i : K_{n_i} \to K_{n_{i + 1}}$ be the composition
$f_{n_{i + 1} - 1} \circ \ldots \circ f_{n_i}$.
We construct commutative diagrams
$$
\vcenter{
\xymatrix{
\bigoplus\nolimits_i K_{n_i} \ar[r]_{1 - g_i} \ar[d]_b &
\bigoplus\nolimits_i K_{n_i} \ar[d]^a \\
\bigoplus\nolimits_n K_n \ar[r]^{1 - f_n} &
\bigoplus\nolimits_n K_n
}
}
\quad\text{and}\quad
\vcenter{
\xymatrix{
\bigoplus\nolimits_n K_n \ar[r]_{1 - f_n} \ar[d]_d &
\bigoplus\nolimits_n K_n \ar[d]^c \\
\bigoplus\nolimits_i K_{n_i} \ar[r]^{1 - g_i} &
\bigoplus\nolimits_i K_{n_i}
}
}
$$
as follows. Let $a_i = a|_{K_{n_i}}$ be the inclusion of $K_{n_i}$
into the direct sum. In other words, $a$ is the natural inclusion.
Let $b_i = b|_{K_{n_i}}$ be the map
$$
K_{n_i}
\xrightarrow{1,\ f_{n_i},\ f_{n_i + 1} \circ f_{n_i},
\ \ldots,\ f_{n_{i + 1} - 2} \circ \ldots \circ f_{n_i}}
K_{n_i} \oplus K_{n_i + 1} \oplus \ldots \oplus K_{n_{i + 1} - 1}
$$
If $n_{i - 1} < j \leq n_i$, then we let $c_j = c|_{K_j}$
be the map
$$
K_j \xrightarrow{f_{n_i - 1} \circ \ldots \circ f_j} K_{n_i}
$$
We let $d_j = d|_{K_j}$ be zero if $j \not = n_i$ for any $i$
and we let $d_{n_i}$ be the natural inclusion of $K_{n_i}$
into the direct sum. In other words, $d$ is the natural projection.
By TR3 these diagrams define morphisms
$$
\varphi : \text{hocolim} K_{n_i} \to \text{hocolim} K_n
\quad\text{and}\quad
\psi : \text{hocolim} K_n \to \text{hocolim} K_{n_i}
$$
Since $c \circ a$ and $d \circ b$ are the identity maps we see that
$\varphi \circ \psi$ is an isomorphism by
Lemma \ref{lemma-third-isomorphism-triangle}.
The other way around we get the morphisms $a \circ c$ and $b \circ d$.
Consider the morphism
$h = (h_j) : \bigoplus K_n \to \bigoplus K_n$ given by
the rule: for $n_{i - 1} < j < n_i$ we set
$$
h_j : K_j
\xrightarrow{1,\ f_j,\ f_{j + 1} \circ f_j,
\ \ldots,\ f_{n_i - 1} \circ \ldots \circ f_j}
K_j \oplus \ldots \oplus K_{n_i}
$$
Then the reader verifies that $(1 - f) \circ h = \text{id} - a \circ c$
and $h \circ (1 - f) = \text{id} - b \circ d$. This means that
$\text{id} - \psi \circ \varphi$ has square zero by
Lemma \ref{lemma-third-map-square-zero} (small argument omitted).
In other words, $\psi \circ \varphi$ differs from the identity
by a nilpotent endomorphism, hence is an isomorphism. Thus
$\varphi$ and $\psi$ are isomorphisms as desired.
\end{proof}

\end{document}
